% ============================================================
% 60_boxes.tex — tcolorbox-Stile, Balken, Boxen, Formelnoten
% ============================================================

\newcommand{\BreakIfNotEnoughSpace}{%
  \par\penalty0\vspace{0pt plus \ZSFbreakThreshold}\penalty-100\vspace{0pt plus -\ZSFbreakThreshold}\relax
}
\newcommand{\BreakIfNotEnoughSpaceChapter}{%
  \par\penalty0\vspace{0pt plus \ZSFbreakThresholdChapter}\penalty-100\vspace{0pt plus -\ZSFbreakThresholdChapter}\relax
}

% ------------------------------------------------------------
% Chapter-Bar
% ------------------------------------------------------------
\newtcolorbox{chapterbar}[1][]{
  enhanced jigsaw,
  breakable=false,
  colback=\ChapterBarColor,
  colframe=\ChapterBarColor,
  boxrule=0pt,
  boxsep=0pt,
  left=\ZSFspaceM,
  right=\ZSFspaceM,
  top=\ZSFspaceM,
  bottom=\ZSFspaceM,
  before skip=16pt,
  after skip=\ZSFspaceM,
  fontupper=\ZSFfontChapter,
  #1
}
\newcommand{\ChapterBar}[1]{
  \BreakIfNotEnoughSpaceChapter
  \penalty-500
  \begin{chapterbar}
    \textcolor{\ChapterBarTextColor}{#1}
  \end{chapterbar}
  \penalty10000
}

% ------------------------------------------------------------
% Dokument-Titel-Masthead (einmalig, vor dem Frontchapter)
% ------------------------------------------------------------
% Kompakter zweireihiger Titelblock in ETH-Blau (ChapterColor0), der
% optisch mit der Frontchapter-Bar zusammenspielt: grosse Zeile =
% Fach/Thema, kleinere kursive Zeile = Autor*in. Einsatz NUR einmal
% ganz oben im Dokument, direkt vor \StartFrontChapter.
\newtcolorbox{zsftitleheader}[1][]{
  enhanced jigsaw,
  breakable=false,
  colback=ZSFDocumentTitleBarColor,
  colframe=ZSFDocumentTitleBarColor,
  boxrule=0pt,
  boxsep=0pt,
  left=\ZSFspaceM,
  right=\ZSFspaceM,
  top=\ZSFspaceS,
  bottom=\ZSFspaceS,
  before skip=0pt,
  after skip=\ZSFspaceXS,
  halign=center,
  #1
}
\newcommand{\ZSFTitleHeader}[2]{%
  \penalty-500
  \begin{zsftitleheader}
    \sffamily
    {\color{ZSFDocumentTitleTextColor}\fontsize{11bp}{12.5bp}\bfseries\selectfont #1\par}%
    \vspace{1pt}%
    {\color{ZSFDocumentTitleBylineColor}\fontsize{6.7bp}{7.5bp}\itshape\selectfont von~#2\par}%
  \end{zsftitleheader}
  \penalty10000
}

% ------------------------------------------------------------
% Subsection-Bar (+ numberless Star-Variante)
% ------------------------------------------------------------
\newcommand{\ZSFSubsectionBarTopSpacing}{%
  \par
  \addvspace{\ZSFsubsectionBarBeforeSkip}%
  \BreakIfNotEnoughSpace
  \ZSFCollapseAfterChapterBar
}
\newtcolorbox{subsectionbar}[1][]{
  enhanced jigsaw,
  breakable=false,
  colback=\SubsectionBarColor,
  colframe=\SubsectionBarColor,
  boxrule=0pt,
  boxsep=0pt,
  left=\ZSFspaceM,
  right=\ZSFspaceM,
  top=\ZSFspaceM,
  bottom=\ZSFspaceM,
  before skip=0pt,
  after skip=\ZSFspaceS,
  fontupper=\ZSFfontSection,
  #1
}
\makeatletter
\newcommand{\SubsectionBar}{\@ifstar{\SubsectionBarStar}{\SubsectionBarNumbered}}
\newcommand{\SubsectionBarNumbered}[2][]{
  \ZSFSubsectionBarTopSpacing
  \penalty-500
  \refstepcounter{subsection}
  \if\relax\detokenize{#1}\relax\else\label{#1}\fi
  \begin{subsectionbar}
    \textcolor{\SubsectionBarTextColor}{\thesubsection. #2}
  \end{subsectionbar}
  \penalty10000
  \ZSFMarkAfterSubsectionBar
}
\newcommand{\SubsectionBarStar}[2][]{
  \ZSFSubsectionBarTopSpacing
  \penalty-500
  \if\relax\detokenize{#1}\relax\else\label{#1}\fi
  \begin{subsectionbar}
    \textcolor{\SubsectionBarTextColor}{#2}
  \end{subsectionbar}
  \penalty10000
  \ZSFMarkAfterSubsectionBar
}
\makeatother

% ------------------------------------------------------------
% Gemeinsamer Basis-Stil für alle Titelbalken (wie Analysis)
% ------------------------------------------------------------
\tcbset{zsftitlebar/.style={
  fonttitle=\ZSFfontBoxTitle\strut,
  halign title=center,
  title style={minimum height=\ZSFtitleBarHeight},
  lefttitle=\ZSFtextMinPad,
  righttitle=\ZSFtextMinPad,
  toptitle=\ZSFspaceXS,
  bottomtitle=\ZSFspaceXS,
  before title={},
  after title={},
}}

% Rechtsbuendiges Meta-/Kategorie-Tag im Titelbalken einer Titelbox.
% Im Titel-Argument verwenden, z.B.:
%   \begin{defbox}[Lorentzkraft\ZSFTitleTag{Magnetostatik}] ... \end{defbox}
\newcommand{\ZSFTitleTag}[1]{\hfill{\ZSFfontTitleTag #1}}

% ------------------------------------------------------------
% Basis-Stil: Panel-Boxen mit Titelbalken (Goal-/Grid-System)
% ------------------------------------------------------------
\tcbset{zsfpanelbox/.style={
  enhanced jigsaw,
  colback=white,
  colbacktitle=\BoxTitleColor,
  colframe=\BoxFrameColor,
  coltitle=\BoxTitleTextColor,
  boxrule=\ZSFtableRuleWidth,
  titlerule=0pt,
  boxsep=0pt,
  arc=\ZSFboxArc,
  outer arc=\ZSFboxArc,
  before skip=\ZSFBoxBeforeSkip,
  after skip=\ZSFspaceS,
  before={\ZSFConsumeAfterChapterBar},
  zsftitlebar,
}}

% ------------------------------------------------------------
% Basis-Stil: Titelboxen (defbox, tablebox, figbox, warnbox)
% ------------------------------------------------------------
% WICHTIG — Single Point of Truth für innere Padding:
% left/right/top/bottom werden HIER im Basis-Stil gesetzt, nicht in
% den einzelnen \newtcolorbox-Definitionen. Damit ist garantiert,
% dass \linewidth innerhalb JEDER Titelbox identisch ist — also auch
% \tabularx{\linewidth}{...} aus styles/20_tables.tex in allen Boxen
% auf dieselbe Breite kommt. Ohne diese Zentralisierung erben
% tablebox/figbox die tcolorbox-Defaults (~4mm) während defbox 4pt
% setzt, was zu unterschiedlichen Tabellen-Breiten führt.
\tcbset{zsftitlebox/.style={
  enhanced jigsaw,
  unbreakable,
  colback=\BoxBackColor,
  colbacktitle=\BoxTitleColor,
  colframe=\BoxFrameColor,
  coltitle=\BoxTitleTextColor,
  boxrule=0.5pt,
  titlerule=0pt,
  boxsep=0pt,
  left=\ZSFspaceM,
  right=\ZSFspaceM,
  top=\ZSFspaceM,
  bottom=\ZSFspaceM,
  before skip=\ZSFBoxBeforeSkip,
  after skip=\ZSFtitleboxAfterSkip,
  before={\ZSFConsumeAfterChapterBar},
  before upper={\parskip=\ZSFspaceS\relax\ZSFReadableAuto},
  zsftitlebar,
}}

\newtcolorbox{defbox}[1][]{zsftitlebox, title={#1}}
% Titelausrichtung der tablebox als eigener Stil: Standard ist zentriert.
% Box-Gruppen, die eine andere Ausrichtung brauchen (Kapitel-Register in
% ch00), definieren den Stil lokal um — siehe ZSFtightBoxGroup.
\tcbset{
  zsftabletitlealign/.style={halign title=center},
  zsftabletitlepadding/.style={toptitle=0pt, bottomtitle=0pt},
  zsftabletitlecolors/.style={colbacktitle=\BoxTitleColor, coltitle=\BoxTitleTextColor},
}

\newtcolorbox{tablebox}[1][]{
  enhanced jigsaw,
  unbreakable,
  colback=white,
  colframe=\ZSFtableFrameColor,
  zsftabletitlecolors,
  boxrule=\ZSFtableFrameWidth,
  titlerule=0pt,
  boxsep=0pt,
  arc=\ZSFboxArc,
  outer arc=\ZSFboxArc,
  clip upper,
  left=0pt,
  right=0pt,
  top=0pt,
  bottom=0pt,
  title={#1},
  fonttitle=\ZSFfontTableBoxTitle\strut,
  zsftabletitlealign,
  title style={minimum height=\ZSFtableTitleBarHeight},
  lefttitle=\ZSFtextMinPad,
  righttitle=\ZSFtextMinPad,
  zsftabletitlepadding,
  before skip=\ZSFBoxBeforeSkip,
  after skip=\ZSFtitleboxAfterSkip,
  before={\ZSFConsumeAfterChapterBar},
  before upper={\parskip=0pt\relax\ZSFReadableAuto},
}

% Titel einer tablebox, die für ein ganzes Kapitel steht (Symbol-Tabellen
% in ch00). Der Titel wird linksbündig gesetzt und trägt rechtsbündig die
% Startseite des Kapitels — die Box-Gruppe wird damit zum Sprungregister.
% Die Seitenzahl kommt aus dem automatischen Kapitel-Label (siehe
% \ZSFMarkChapterPage in styles/40_colors_structure.tex).
% Voraussetzung ist die linksbündige Titelausrichtung des umgebenden
% \ZSFtightBoxGroup (siehe zsftabletitlealign) — erst dadurch schiebt das
% \hfill die Seitenzahl an den rechten Rand.
% Beide Argumente sind Pflicht (kein optionales): der Titel steht im
% optionalen Argument von \begin{tablebox}[...], eine geschachtelte
% []-Klammer würde dieses vorzeitig beenden.
% Argumente: {<Kapitelnummer>}{<Titel>}
\newcommand{\ZSFchapterTableTitle}[2]{#1.~#2\hfill(S.~\ZSFChapterPage{#1})}

% ------------------------------------------------------------
% Legenden-Band (gemeinsamer Spaltenkopf über einer Box-Gruppe)
% ------------------------------------------------------------
% Wenn mehrere tablebox-Blöcke dieselbe Spaltenstruktur haben, ist eine
% Kopfzeile pro Box redundant. \begin{ZSFlegendband}[<font>]{<colspec>}
% rendert den Kopf EINMAL als vollbreites Farbband über der Gruppe.
%
% Geometrie bewusst an tablebox gekoppelt: tablebox setzt left/right=0pt
% bei boxrule=\ZSFsubtleRuleWidth, der Inhalt beginnt also um genau die
% Rahmenstärke eingerückt. Das Band spiegelt das (boxrule=0pt,
% left/right=\ZSFsubtleRuleWidth), damit die Bandspalten exakt über den
% Tabellenspalten der Boxen sitzen. Wird eine der beiden Geometrien
% geändert, muss die andere mitgezogen werden.
\newtcolorbox{ZSFlegendbandFrame}{
  enhanced jigsaw,
  unbreakable,
  colback=\ZSFheaderRowBG,
  colframe=\ZSFheaderRowBG,
  boxrule=0pt,
  boxsep=0pt,
  arc=\ZSFboxArc,
  outer arc=\ZSFboxArc,
  left=\ZSFsubtleRuleWidth,
  right=\ZSFsubtleRuleWidth,
  top=0pt,
  bottom=0pt,
  before skip=\ZSFBoxBeforeSkip,
  after skip=\ZSFtitleboxAfterSkip,
  before={\ZSFConsumeAfterChapterBar},
  before upper={\parskip=0pt\relax\ZSFReadableAuto},
}
% Bewusst ein Makro und keine Umgebung: tcolorbox setzt den Inhalt über eine
% interne savebox, die zwingend über \begin/\end laufen muss — als geteilte
% Umgebung (Begin-/End-Code) kollidiert das mit dem \@currenvir-Check.
% Argumente: [<font>]{<colspec>}{<Kopfzeile inkl. \\>}
\NewDocumentCommand{\ZSFlegendband}{O{\ZSFfontTable} m m}{%
  \begin{ZSFlegendbandFrame}%
    \begin{ZSFtablePlain}[#1]{#2}%
      #3%
    \end{ZSFtablePlain}%
  \end{ZSFlegendbandFrame}%
}
% figbox: Grafik-Titelbox. Top/Bottom werden gegenüber defbox/tablebox
% reduziert (\ZSFspaceS statt \ZSFspaceM) — Bilder bringen visuell
% schon eigene Ränder mit; weniger Innenpadding spart Platz.
\newtcolorbox{figbox}[1][]{%
  zsftitlebox,
  title={#1},
  toptitle=\ZSFspaceS,
  top=\ZSFspaceS,
  bottom=\ZSFspaceS,
}

% ------------------------------------------------------------
% Formula-Box + Longformula
% ------------------------------------------------------------
\newenvironment{longformula}{\[\begin{aligned}}{\end{aligned}\]}
\newtcolorbox{formulabox}[1][]{
  enhanced jigsaw,
  unbreakable,
  colback=\FormulaboxBackColor,
  colframe=black,
  colbacktitle=\BoxTitleColor,
  coltitle=\BoxTitleTextColor,
  boxrule=0.5pt,
  boxsep=0pt,
  left=\ZSFspaceM,
  right=\ZSFspaceM,
  top=\ZSFspaceM,
  bottom=\ZSFspaceM,
  before skip=\ZSFspaceS,
  after skip=\ZSFspaceS,
  before={\ZSFConsumeAfterChapterBar},
  before upper={\parskip=\ZSFspaceS\centering},
  after upper={\par},
  middle=\ZSFspaceSep,
  title={#1},
  zsftitlebar,
}

% Dünne, graue Trennlinie für gestapelte Formeln in einer formulabox.
% Einsatz: zwischen zwei \[ ... \] innerhalb derselben formulabox.
\newcommand{\formulasep}{%
  \par\vspace{-\parskip}\vspace{\ZSFspaceSep}%
  \noindent\textcolor{black!30}{\rule{\linewidth}{0.5pt}}%
  \par\vspace{-\parskip}\vspace{\ZSFspaceSep}%
  \penalty-50%
}

% Zentriertes Item-Heading mit Trennstrich — für Listen gleichwertiger
% Teilaspekte (z. B. Fälle, Varianten) innerhalb einer Box.
\newcommand{\ZSFItemHeading}[1]{%
  \par\vspace{-\parskip}\vspace{\ZSFspaceM}%
  {\centering\textbf{#1}\par}\vspace{-\parskip}\vspace{\ZSFspaceS}%
  \noindent\rule{\linewidth}{0.5pt}\par\vspace{-\parskip}\vspace{\ZSFspaceS}%
}

\newcommand{\formulanote}[1]{%
  \par\vspace{-\parskip}\vspace{\ZSFspaceS}%
  \noindent{{\ZSFfontNote #1}}%
  \par\vspace{-\parskip}\vspace{\ZSFspaceS}%
}
\newcommand{\formulanoteabove}[1]{%
  \par\vspace{-\parskip}\vspace{\ZSFspaceS}%
  \noindent{{\ZSFfontNote #1}}%
  \par\vspace{-\parskip}\vspace{\ZSFspaceS}%
}

% ------------------------------------------------------------
% Warn-Box (rot, Physik-spezifisch)
% ------------------------------------------------------------
\newtcolorbox{warnbox}[1][Achtung]{
  zsftitlebox,
  colback=WarnboxBack,
  colbacktitle=ETHRot,
  coltitle=white,
  title=#1,
}

% ------------------------------------------------------------
% Runintext + Grafik-Platzhalter
% ------------------------------------------------------------
\newenvironment{runintext}{\ZSFConsumeAfterChapterBar\par\addvspace{\ZSFspaceS}\noindent\ZSFReadableAuto\ignorespaces}{\par}

\newcommand{\PlaceholderGraphic}[1]{
  \fbox{\parbox{0.9\linewidth}{\centering\small [Grafik: #1]}}
}

% ------------------------------------------------------------
% Grafik-Helpers — einheitliche Einbindung von Bildern in figboxen.
% ------------------------------------------------------------
% \ZSFfigcaption{...} — dezente Bildunterschrift unter einem Bild.
% Leerer Text wird komplett weggelassen (kein Zusatzabstand).
% Style: \ZSFfontNote (sffamily italic 6.7bp) + dezenter Grauton.
% Manuelle Nutzung möglich, wird aber i.d.R. von \ZSFfig gerufen.
\newcommand{\ZSFfigcaption}[1]{%
  \if\relax\detokenize{#1}\relax\else
    \par\vspace{\ZSFspaceXS}%
    {\ZSFfontNote\color{black!60}#1\par}%
  \fi
}

% \ZSFfig{graphics/foo.png}{Titel}{Kurzbeschreibung}
% Ein-Zeilen-Ersatz für das wiederkehrende Muster
%   \begin{figbox}[Titel]\centering\includegraphics[...]{...}\end{figbox}
% Vorteile:
%   * Einheitliche Max-Grösse via \ZSFfigMaxHeight (verhindert, dass
%     hohe Bilder "overkill" Platz einnehmen).
%   * Optionale Bildunterschrift über den 3. Parameter.
%     Leere Beschreibung ({}) → keine Caption, kein Zusatzabstand.
% Hinweis: Titel mit Kommata oder Math-Mode bitte in {...} wrappen,
%   z.B. \ZSFfig{pfad}{{Sinus U(t), I(t)}}{...}
\newcommand{\ZSFfig}[4][\ZSFfigMaxHeight]{%
  \begin{figbox}[{#3}]%
    \centering
    \includegraphics[width=\linewidth,height=#1,keepaspectratio]{#2}%
    \ZSFfigcaption{#4}%
  \end{figbox}%
}

% \ZSFfigside{graphics/foo.png}{Titel}{Rechte Seite}
% Zweispalten-Variante zu \ZSFfig: Bild links (ca. 30%), Inhalt rechts
% (ca. 66%), beide vertikal zentriert. Für platzsparende Reihen
% ähnlicher Objekte (z.B. Kondensator-Geometrien, ebene Kurven).
%
% Der rechte Content darf frei strukturiert werden:
%   * Formeln als \[ ... \]
%   * Bold-Labels via \textVorFormel{Kapazität:}
%   * Trennlinie zwischen Sektionen via \formulasep
%   * Kleine Fussnote via \formulanote{...}
%   * Freier Fliesstext (\small für kompakt)
%
% Beispiel:
%   \ZSFfigside{graphics/bar.png}{Plattenkondensator}{%
%     \textVorFormel{E-Feld:}
%     \[ E = U/d \]
%     \formulasep
%     \textVorFormel{Kapazität:}
%     \[ C = \varepsilon_0 A/d \]
%   }
% Die Breiten werden ERST im Kontext der Box ausgewertet — \linewidth ist in
% der Präambel die volle Textbreite (alle Spalten), nicht die Spaltenbreite.
\newlength{\ZSFfigsideImgWidth}
\newlength{\ZSFfigsideTextWidth}
\newcommand{\ZSFfigside}[3]{%
  \begin{figbox}[{#2}]%
    \noindent
    \setlength{\ZSFfigsideImgWidth}{0.32\linewidth}%
    \setlength{\ZSFfigsideTextWidth}{0.64\linewidth}%
    \begin{minipage}[c]{\ZSFfigsideImgWidth}%
      \centering
      \includegraphics[width=\linewidth,height=\ZSFfigMaxHeight,keepaspectratio]{#1}%
    \end{minipage}%
    \hfill
    \begin{minipage}[c]{\ZSFfigsideTextWidth}%
      \raggedright
      #3%
    \end{minipage}%
    \par
  \end{figbox}%
}

% --- Nebenläufiges Bild/Text-Layout (für "Ebene Kurven", Tabellen mit Bildern etc.) ---
\newtcolorbox{splitbox}[1][0.3]{
  blankest,
  sidebyside,
  sidebyside align=top,
  sidebyside gap=\ZSFspaceM,
  lefthand width=#1\linewidth,
  lower separated=false,
  before={\ZSFConsumeAfterChapterBar},
  before skip=\ZSFspaceS,
  after skip=\ZSFspaceS,
  before upper={\ZSFReadableAuto},
  before lower={\ZSFReadableAuto},
  boxsep=0pt, left=0pt, right=0pt, top=0pt, bottom=0pt
}

% ------------------------------------------------------------
% Aussagen / Verfahren / Listen
% ------------------------------------------------------------
% statementbox: dezente Aussagebox mit linkem Farbbalken.
\tcbset{zsfstatementstyle/.style={
  enhanced jigsaw,
  breakable=false,
  colback=white,
  colframe=white,
  frame hidden,
  borderline west={\ZSFstatementBarWidth}{0pt}{\StatementBarColor},
  boxrule=0pt,
  arc=0pt,
  outer arc=0pt,
  left=\ZSFstatementLeftPad,
  right=\ZSFspaceS,
  top=\ZSFspaceXS,
  bottom=\ZSFspaceXS,
  before skip=\ZSFBoxBeforeSkip,
  after skip=\ZSFspaceS,
  before={\ZSFConsumeAfterChapterBar},
  before upper={\ZSFReadableAuto},
}}

\NewDocumentEnvironment{statementbox}{O{}}
  {\begin{tcolorbox}[zsfstatementstyle]%
   \if\relax\detokenize{#1}\relax\else
     {\ZSFfontStatementTitle\color{\StatementTitleColor}#1}\par\vspace{\ZSFspaceXS}%
   \fi\ignorespaces}
  {\end{tcolorbox}}

% procedure: nummerierte Schritt-Liste in einer statementbox.
% Verwendung: \ProcStep{Titel}{Beschreibung}
\NewDocumentEnvironment{procedure}{O{}}
  {\begin{statementbox}[#1]%
   \begin{enumerate}[
     leftmargin=\ZSFprocLeftMargin,
     labelsep=\ZSFprocLabelSep,
     itemsep=\ZSFprocItemSep,
     topsep=0pt,
     parsep=0pt,
     partopsep=0pt
   ]}
  {\end{enumerate}\end{statementbox}}
\newcommand{\ProcStep}[2]{\item {\ZSFfontProcedureStep #1}\\[\ZSFprocStepGap]#2}

% factlist: rahmenlose Liste von Definitions-/Eigenschaftsfakten.
% Verwendung: \ZSFFact{Lead-in}{Beschreibung}
\NewDocumentEnvironment{factlist}{}
  {\begin{itemize}[
     leftmargin=\ZSFfactLeftMargin,
     labelsep=\ZSFfactLabelSep,
     label={\textcolor{\StatementBarColor}{\textbullet}},
     itemsep=\ZSFprocItemSep,
     topsep=\ZSFspaceXS,
     parsep=0pt,
     partopsep=0pt
   ]}
  {\end{itemize}}
\newcommand{\ZSFFact}[2]{\item {\ZSFfontFactLead #1}\,---\,#2}

% propertylist: factlist innerhalb einer statementbox.
\NewDocumentEnvironment{propertylist}{O{}}
  {\begin{statementbox}[#1]\begin{itemize}[
     leftmargin=\ZSFfactLeftMargin,
     labelsep=\ZSFfactLabelSep,
     itemsep=\ZSFprocItemSep,
     topsep=0pt,
     parsep=0pt,
     partopsep=0pt
   ]}
  {\end{itemize}\end{statementbox}}

% ------------------------------------------------------------
% Goal-System: Bedingung -> Umformung -> Ziel
% ------------------------------------------------------------
\newtcolorbox{goalbox}[1][]{
  zsfpanelbox,
  breakable=false,
  colback=white,
  left=\ZSFspaceS,
  right=\ZSFspaceS,
  top=\ZSFtextMinPad,
  bottom=\ZSFtextMinPad,
  title={#1},
  before upper={\parskip=\ZSFspaceS\centering},
  after upper={\par},
}

\newtcbox{\GoalConditionBox}{%
  on line,
  enhanced,
  boxrule=\ZSFemphasisRuleWidth,
  colback=\GoalConditionBackColor,
  colframe=\GoalConditionFrameColor,
  borderline={\ZSFsubtleRuleWidth}{0pt}{\GoalConditionBorderColor},
  arc=\ZSFboxArc,
  left=\ZSFspaceXS,
  right=\ZSFspaceXS,
  top=\ZSFspaceXS,
  bottom=\ZSFspaceXS,
  nobeforeafter,
  tcbox raise base,
}

\newtcbox{\GoalTargetBox}{%
  on line,
  enhanced,
  boxrule=\ZSFtableRuleWidth,
  colback=\GoalTargetBackColor,
  colframe=\GoalTargetFrameColor,
  arc=\ZSFboxArc,
  left=\ZSFspaceXS,
  right=\ZSFspaceXS,
  top=\ZSFspaceXS,
  bottom=\ZSFspaceXS,
  nobeforeafter,
  tcbox raise base,
}

\newcommand{\GoalCondition}[1]{\par\centering\GoalConditionBox{\ZSFfontNote #1}\par}
\newcommand{\GoalStep}[1]{\par\centering$\displaystyle #1$\par}
\newcommand{\GoalTarget}[1]{\par\centering\GoalTargetBox{$\displaystyle #1$}\par}
\newcommand{\ZSFDerivationCase}[3]{%
  \GoalCondition{#1}%
  \GoalStep{#2 =}%
  \GoalTarget{#3}%
}
\newcommand{\ZSFDerivationInlineCase}[3]{%
  \GoalCondition{#1}%
  \GoalStep{#2 = \GoalTargetBox{$\displaystyle #3$}}%
}
\newcommand{\GoalCard}[4][]{%
  \begin{goalbox}[#1]
    \GoalCondition{#2}
    #3
    \GoalTarget{#4}
  \end{goalbox}%
}
\newcommand{\GoalPair}[8]{%
  \begin{tcbraster}[raster columns=2, raster column skip=\ZSFspaceM, raster row skip=0pt, raster equal height]
    \GoalCard[#1]{#2}{#3}{#4}
    \GoalCard[#5]{#6}{#7}{#8}
  \end{tcbraster}%
}

% Inline-Danger-Tag für Stolperfallen im Fliesstext.
\newtcbox{\ZSFdanger}{
  on line,
  colback=DangerBackColor,
  colframe=DangerBackColor,
  coltext=DangerTextColor,
  boxsep=\ZSFspaceXS,
  left=\ZSFspaceS,
  right=\ZSFspaceS,
  top=\ZSFspaceXS,
  bottom=\ZSFspaceXS,
  boxrule=0pt,
  arc=\ZSFboxArc,
  fontupper=\ZSFfontDanger,
}

% ------------------------------------------------------------
% Valuegrid: kompakte tabularray-Wertegrids in einer Panel-Box
% ------------------------------------------------------------
\newtcolorbox{compactgridbox}[1][]{
  zsfpanelbox,
  colframe=\CompactGridFrameColor,
  breakable=false,
  title={#1},
  boxsep=0pt,
  left=0pt,
  right=0pt,
  top=\ZSFtextMinPad,
  bottom=\ZSFtextMinPad,
  boxrule=\ZSFtableRuleWidth,
}
\newcommand{\ZSFvaluegridCommon}{
  colsep=\ZSFvaluegridColSep,
  rowsep=\ZSFvaluegridRowSep,
  hlines,
  vlines,
  rows={valign=m}
}
\NewDocumentEnvironment{valuegrid}{m O{}}
  {\begin{compactgridbox}[#2]\begin{tblr}{colspec={l *{#1}{c}},\ZSFvaluegridCommon}}
  {\end{tblr}\end{compactgridbox}}
\NewDocumentEnvironment{valuegridcustom}{m O{}}
  {\begin{compactgridbox}[#2]\begin{tblr}{colspec={#1},\ZSFvaluegridCommon}}
  {\end{tblr}\end{compactgridbox}}

% Legacy-Aliase für fixe Datenspaltenzahlen.
\NewDocumentEnvironment{valuegridtwo}{O{}}{\begin{valuegrid}{2}[#1]}{\end{valuegrid}}
\NewDocumentEnvironment{valuegridthree}{O{}}{\begin{valuegrid}{3}[#1]}{\end{valuegrid}}
\NewDocumentEnvironment{valuegridfour}{O{}}{\begin{valuegrid}{4}[#1]}{\end{valuegrid}}
\NewDocumentEnvironment{valuegridfive}{O{}}{\begin{valuegrid}{5}[#1]}{\end{valuegrid}}
\NewDocumentEnvironment{valuegridsix}{O{}}{\begin{valuegrid}{6}[#1]}{\end{valuegrid}}
\NewDocumentEnvironment{valuegridseven}{O{}}{\begin{valuegrid}{7}[#1]}{\end{valuegrid}}
